%═══════════════════════════════════════════════════════════════════════════════
\section{More Robustness: The Doubly-Infinite Tape (Exercise 2, Tutorial 2)} \kemark
%═══════════════════════════════════════════════════════════════════════════════

This section covers \textbf{Exercise 2 from Tutorial 2}. The exercise has three
parts, and we walk through each one in full detail. By the end, you will have a
fourth robustness result: changing the shape of the tape does not add
computational power.

%───────────────────────────────────────────────────────────────────────────────
\subsection{What Is a Doubly-Infinite Tape?}
\label{subsec:doubly-infinite-what}
%───────────────────────────────────────────────────────────────────────────────

\begin{notebox}{Analogy: Receipt Paper vs.\ a Strip Floating in Space}
Picture a receipt printer at a grocery store. The paper roll has a clear
\textbf{beginning} --- the point where the paper is attached to the roll ---
and it unrolls infinitely to one side (you can always print more). If you try
to pull the paper \emph{backwards} past the attachment point, you can't ---
there's nothing there. That beginning is a hard boundary.

Now imagine you rip the receipt off the roll and let it float in zero gravity
in the middle of a room. There is no ``start'' and no ``end'' anymore. You
can drift to the left as far as you want, or drift to the right as far as you
want. The strip extends infinitely in \emph{both} directions.

\textbf{In our formal world:}
\begin{itemize}[nosep]
\item Receipt paper roll = \textbf{singly-infinite tape} (standard TM). The
  attachment point is cell 1 --- the leftmost cell. Cells are numbered
  $1, 2, 3, \ldots$ The head cannot go left of cell 1.
\item Strip floating in space = \textbf{doubly-infinite tape}. No left wall,
  no right wall. Cells are numbered $\ldots, -2, -1, 0, 1, 2, \ldots$ (all
  integers $\mathbb{Z}$). The head can go left forever.
\item Attachment point = position 1 boundary on a singly-infinite tape.
\item No attachment point = no boundary at all on a doubly-infinite tape.
\end{itemize}
\end{notebox}

\begin{conceptbox}{Why Consider a Doubly-Infinite Tape?}
\textbf{The Core Problem:}

On a standard (singly-infinite) TM, when the head is at position 1 and the
transition says ``move left,'' the head \textbf{stays at position 1}. It gets
``stuck'' against the left wall. This means:
\begin{itemize}[nosep]
\item Every TM algorithm must \emph{remember} that the left boundary exists.
\item Some algorithms require awkward workarounds to avoid running into the
  wall (e.g., placing a special marker at the leftmost cell so you know
  you've reached the edge).
\end{itemize}

On a doubly-infinite tape, you can \textbf{always go left}. The tape extends
infinitely in both directions. Some algorithms become simpler to design when
you don't have to worry about the left boundary.

\textbf{The Question:} Does this convenience add \textbf{computational power}?
Can a doubly-infinite tape TM decide languages that a standard TM cannot?

Spoiler: \textbf{No.} The two models are equivalent. That is exactly what
Exercise 2 asks you to prove, and that is exactly what we will trace through
in this section.

\kt{The doubly-infinite tape removes the left boundary of a standard TM.
The question is whether this makes the machine more powerful --- and the
answer is no. This is our fourth robustness result.}
\end{conceptbox}

%───────────────────────────────────────────────────────────────────────────────
\subsection{Part 1: Formal Definition of a Doubly-Infinite Tape DTM}
\label{subsec:doubly-infinite-formal}
%───────────────────────────────────────────────────────────────────────────────

\begin{notebox}{{Analogy: Same Vending Machine, Different Shelf Layout}}
Think of a vending machine. The \textbf{machine itself} --- its buttons, its
coin slot, its dispensing mechanism --- doesn't change. What changes is the
\textbf{shelf layout inside}: instead of shelves labeled $1, 2, 3, \ldots$
(starting from shelf 1 at the top), you now have shelves labeled
$\ldots, -2, -1, 0, 1, 2, \ldots$ (extending upward and downward
indefinitely). The robot arm that fetches items can now reach shelves in
\emph{both} directions.

\textbf{In our formal world:}
\begin{itemize}[nosep]
\item Vending machine = the TM tuple $(Q, \Sigma, \Gamma, s, t, r, \delta)$
  --- this does NOT change.
\item Shelf layout = the tape and how cells are indexed --- this is what
  changes (from $\mathbb{N}$ to $\mathbb{Z}$).
\item Robot arm = the read/write head --- now it can reach negative-numbered
  cells too.
\end{itemize}
\end{notebox}

\begin{definitionbox}{{Doubly-Infinite Tape DTM --- The Tuple}}
\textbf{Plain English:}
A doubly-infinite tape DTM has \textbf{exactly the same 7-tuple} as a standard
DTM. The states, alphabets, start/accept/reject states, and transition
function are all identical. What changes is not the machine's description but
\textbf{how the machine operates} --- specifically, how we define
configurations and how the head moves.

\textbf{Formal:}
A doubly-infinite tape DTM is a 7-tuple:
\fitmath{M = (Q,\; \Sigma,\; \Gamma,\; s,\; t,\; r,\; \delta)}
with exactly the same meaning as for a standard DTM:
\begin{itemize}[nosep]
\item $Q$ = finite set of states
\item $\Sigma$ = input alphabet (no blank symbol)
\item $\Gamma \supseteq \Sigma \cup \{\sqcup\}$ = tape alphabet
\item $s \in Q$ = start state
\item $t \in Q$ = accept state
\item $r \in Q$, $r \neq t$ = reject state
\item $\delta \colon (Q \setminus \{t,r\}) \times \Gamma \to Q \times \Gamma \times \{-1,+1\}$ = transition function
\end{itemize}

The transition function is identical: read a symbol, write a symbol, move left
($-1$) or right ($+1$). Nothing in the tuple changes.

\textbf{What changes is the SEMANTICS} --- the meaning of ``configuration''
and ``successor configuration.'' That's in the next box.
\end{definitionbox}

\begin{definitionbox}{{Doubly-Infinite Tape DTM --- Configurations}}
\textbf{Plain English:}
A configuration is a snapshot of the machine at one moment in time: what state
it's in, what's on the tape, and where the head is. For a doubly-infinite
tape, the tape content is a function from \emph{all integers} (not just
positive integers) to tape symbols. The head position can be any integer,
including negative numbers.

\textbf{Formal:}
A \textbf{configuration} is a triple $[q, \tau, \ell]$ where:
\begin{itemize}[nosep]
\item $q \in Q$ is the current state
\item $\tau \colon \mathbb{Z} \to \Gamma$ is the tape content
\item $\ell \in \mathbb{Z}$ is the head position
\end{itemize}

\textbf{Breaking Down the Notation:}
\begin{itemize}[nosep]
\item $\mathbb{Z}$ means \textbf{all integers}: $\ldots, -3, -2, -1, 0, 1, 2, 3, \ldots$ This is the key difference from a standard TM, where the tape is indexed by $\mathbb{N} = \{1, 2, 3, \ldots\}$.
\item $\tau \colon \mathbb{Z} \to \Gamma$ means every integer position maps to some tape symbol. Position $-5$ has a symbol, position $0$ has a symbol, position $42$ has a symbol.
\item $\ell \in \mathbb{Z}$ means the head can be at position $-3$, or $0$, or $17$ --- any integer at all, including negative ones. On a standard TM, the head position is always $\geq 1$.
\end{itemize}

\textbf{Initial configuration} on input $w = w_1 w_2 \cdots w_k$:
\fitmath{[s,\; \tau_0,\; 1]}
where:
\fitmath{\tau_0(n) = \begin{cases} w_n & \text{if } n \in \{1, \ldots, k\} \\ \sqcup & \text{otherwise (including } n = 0, -1, -2, \ldots\text{)} \end{cases}}

\textbf{In words:} The input is placed at positions $1$ through $k$ (just
like a standard TM). Every other cell --- including all the negative
positions and zero --- contains the blank symbol $\sqcup$. The head starts at
position 1, in the start state $s$.
\end{definitionbox}

\begin{definitionbox}{{Doubly-Infinite Tape DTM --- Successor Configuration}}
\textbf{Plain English:}
Given a configuration $[q, \tau, \ell]$, the machine reads the symbol at the
head position, looks up the transition, writes a new symbol, changes state,
and moves the head. This is exactly the same as a standard TM --- with one
critical difference: there is \textbf{no special case} for the head being at
the leftmost cell.

\textbf{Formal:}
If $\delta(q, \tau(\ell)) = (p, a, d)$ where $d \in \{-1, +1\}$, then the
successor configuration is:
\fitmath{[p,\; \tau',\; \ell + d]}
where:
\fitmath{\tau'(n) = \begin{cases} a & \text{if } n = \ell \\ \tau(n) & \text{otherwise} \end{cases}}

\textbf{Breaking Down the Notation:}
\begin{itemize}[nosep]
\item $\tau(\ell)$ = the symbol currently under the head (at position $\ell$)
\item $(p, a, d)$ = the transition's output: go to state $p$, write symbol
  $a$, move in direction $d$
\item $\tau'$ = the updated tape: identical to $\tau$ everywhere except at
  position $\ell$, where the old symbol is replaced by $a$
\item $\ell + d$ = the new head position. If $d = +1$, head moves right. If
  $d = -1$, head moves left.
\end{itemize}

\textbf{THE KEY DIFFERENCE from a standard TM:}

On a standard singly-infinite tape, the successor rule has a special case:
\begin{quote}
``If $\ell = 1$ and $d = -1$, the head stays at position 1 (cannot go left
of the boundary).''
\end{quote}

On a doubly-infinite tape, \textbf{this special case does not exist}. If the
head is at position 1 and the transition says ``move left,'' the head goes to
position \textbf{0}. If the head is at position 0 and says ``move left,'' it
goes to position $-1$. The head can always go left, without limit.

\textbf{Accepting, rejecting, and halting:} Same as standard DTM. The machine
accepts when it enters state $t$, rejects when it enters state $r$, and halts
when it enters either.
\end{definitionbox}

\begin{examplebox}{Concrete Example: Input ``01'' on a Doubly-Infinite Tape}
Let's see exactly what the tape looks like for input $w = 01$ (so $k = 2$,
$w_1 = 0$, $w_2 = 1$).

\textbf{Initial configuration:} $[s, \tau_0, 1]$

The tape content $\tau_0$ assigns symbols to every integer:
\begin{center}
\begin{tabular}{c|cccccccc}
Position & $\cdots$ & $-2$ & $-1$ & $0$ & $1$ & $2$ & $3$ & $\cdots$ \\
\hline
Symbol & $\cdots$ & $\sqcup$ & $\sqcup$ & $\sqcup$ & $\mathtt{0}$ & $\mathtt{1}$ & $\sqcup$ & $\cdots$ \\
\end{tabular}
\end{center}

The head is at position 1 (pointing at the $\mathtt{0}$). All positions outside
$\{1, 2\}$ contain blanks --- including position 0, $-1$, $-2$, and so on.

\textbf{What happens if the head moves left from position 1?}

On a \textbf{standard TM}: the head stays at position 1 (stuck against the
left wall).

On a \textbf{doubly-infinite tape TM}: the head moves to position 0. It reads
$\sqcup$ (blank) there. If it moves left again, it goes to position $-1$ ---
also blank. There is no wall. The head just keeps finding blank cells as it
moves further left.

\kt{The only difference between a singly-infinite and doubly-infinite tape DTM
is in the successor configuration: no left boundary means the head can always
move left, finding blank cells extending infinitely to the left.}
\end{examplebox}


%───────────────────────────────────────────────────────────────────────────────
\subsection{Part 2: Singly-Infinite Simulates Doubly-Infinite (The Hard Direction)}
\label{subsec:singly-simulates-doubly}
%───────────────────────────────────────────────────────────────────────────────

This is the hard direction. A singly-infinite tape has positions $1, 2, 3,
\ldots$ but a doubly-infinite tape needs $\ldots, -2, -1, 0, 1, 2, \ldots$
How do you fit an infinite-in-both-directions tape onto a tape that only goes
one way?

\begin{notebox}{Analogy: Fitting a Two-Way Street onto a One-Way Street}
Imagine you have a street that goes in \textbf{both directions} --- cars drive
east and west. Now the city needs to convert it into a \textbf{one-way street}
(only east). But you still need to handle all the traffic that used to go
west.

The solution: put up \textbf{barriers} at the beginning and end of the active
zone. Whenever a car needs to go west past the barrier, you \textbf{shift} all
the parked cars one space to the right to make room, then let the car park in
the newly created space. The one-way street handles the same traffic --- it
just requires occasional reshuffling.

\textbf{In our formal world:}
\begin{itemize}[nosep]
\item Two-way street = doubly-infinite tape (extends left and right)
\item One-way street = singly-infinite tape (extends only right)
\item Traffic going west = the head trying to move to negative positions
\item Barriers = special boundary markers $\ell$ (left) and $r$ (right)
\item Shifting parked cars = shifting tape content one cell right to create
  space for a new blank cell on the left
\end{itemize}
\end{notebox}

\begin{conceptbox}{The Big Idea of the Construction}
\textbf{The Core Insight:}

We cannot physically add negative-numbered cells to a singly-infinite tape.
But we can maintain a ``window'' of cells that contains everything the
doubly-infinite machine has ever touched. This window is marked by two special
symbols:
\begin{itemize}[nosep]
\item $\ell$ = left boundary marker (``the leftmost cell $M$ has visited'')
\item $r$ = right boundary marker (``the rightmost cell $M$ has visited'')
\end{itemize}

The content between $\ell$ and $r$ is a faithful copy of the corresponding
region of $M$'s doubly-infinite tape. Whenever $M$ tries to move \emph{outside}
this window (to a cell it hasn't visited before), our simulator \textbf{expands
the window} to accommodate it.

\textbf{Why This Matters:}

At any finite point in a computation, $M$ has only visited finitely many cells.
Even though the doubly-infinite tape is infinite in both directions, the
machine has only used a finite portion. Our window captures exactly that
portion.

\kt{The key insight: a doubly-infinite tape machine only uses finitely many
cells at any finite step. We track those cells with boundary markers and
expand as needed.}
\end{conceptbox}

\begin{definitionbox}{{The Simulation Construction ($M'$ Simulates $M$)}}
\textbf{Plain English:}
Given a doubly-infinite tape DTM $M$, we build a singly-infinite tape DTM
$M'$ that outcome-simulates $M$. The simulator $M'$ uses two special markers
not in $M$'s tape alphabet to keep track of boundaries, and it ``expands'' the
tracked region whenever $M$ needs a new cell.

\textbf{The construction:}

Let $M = (Q, \Sigma, \Gamma, s, t, r, \delta)$ be a doubly-infinite tape DTM.
We build $M'$ with tape alphabet
$\Gamma' = \Gamma \cup \{\ell, r\}$ where $\ell, r \notin \Gamma$ are fresh
symbols (the boundary markers).

$M'$ operates as follows:

\textbf{Step (a) --- Initialization:}
Replace the input $w = w_1 \cdots w_k$ on the tape by
$\ell\, w_1 \cdots w_k\, r$. That is, place marker $\ell$ immediately before
the input and marker $r$ immediately after. Position the head at the first
letter of $w$ (or at $r$ if $w = \varepsilon$).

\begin{lstlisting}
Before: [ w1 | w2 | ... | wk | _ | _ | ... ]
After:  [  l | w1 | w2 | ... | wk | r | _ | _ | ... ]
                ^
          head here (at w1)
\end{lstlisting}

\textbf{Step (b) --- Check for halt:}
If the current simulated state is $t$ (accept) $\to$ accept.
If the current simulated state is $r_{\text{rej}}$ (reject) $\to$ reject.

\textbf{Step (c) --- Simulate one transition of $M$:}
Read the symbol under the head (ignoring $\ell$ and $r$ markers --- the head
should never be reading a marker during normal simulation). Look up
$\delta(q, a)$ where $q$ is the current simulated state and $a$ is the symbol
read. Write the new symbol, update the simulated state, and move the head
according to $\delta$.

\textbf{Step (d) --- Left expansion (head hit $\ell$):}
If after step (c) the head is pointing at the $\ell$ marker, that means $M$
wants to access a cell to the left of the current window. We need to create
space:
\begin{enumerate}[nosep]
\item Shift ALL content from $\ell$ to $r$ (inclusive) one cell to the right.
  This frees up one cell immediately after $\ell$.
\item Write $\sqcup$ (blank) in the freed cell --- this represents the new
  cell $M$ is visiting for the first time (which would be blank on $M$'s tape).
\item Move the head to that blank cell.
\end{enumerate}

\begin{lstlisting}
Before shift: [ l | A | B | C | r | _ | ... ]  head at pos 1 (l)
After shift:  [ l | _ | A | B | C | r | _ | ... ]  head at pos 2 (blank)
\end{lstlisting}

\textbf{Step (e) --- Right expansion (head hit $r$):}
If after step (c) the head is pointing at the $r$ marker, $M$ wants to access
a cell to the right of the current window. This is simpler:
\begin{enumerate}[nosep]
\item Replace $r$ with $\sqcup$ (the new cell, which would be blank on
  $M$'s tape).
\item Write $r$ one cell further to the right (extending the right boundary).
\item The head stays on the newly created blank cell.
\end{enumerate}

\begin{lstlisting}
Before: [ l | A | B | C | r | _ | ... ]  head at pos 5 (r)
After:  [ l | A | B | C | _ | r | ... ]  head at pos 5 (now blank)
\end{lstlisting}

\textbf{Step (f):} Go back to step (b).

\textbf{Breaking Down the Key Ideas:}
\begin{itemize}[nosep]
\item \textbf{$\ell$ and $r$ markers:} These are NOT part of $M$'s alphabet.
  They are bookkeeping symbols used only by $M'$ to know where the
  ``active window'' is. $M$ never sees them.
\item \textbf{Left expansion vs.\ right expansion:} Left expansion requires
  shifting (expensive --- $O(n)$ work where $n$ is the window size). Right
  expansion just moves the $r$ marker (cheap --- $O(1)$ work). This asymmetry
  exists because the singly-infinite tape extends freely to the right but has
  a wall on the left.
\item \textbf{Blank cells:} When $M$ moves to a cell it has never visited
  before, that cell contains $\sqcup$ on $M$'s tape. Our expansion steps
  ensure that the corresponding cell on $M'$'s tape also contains $\sqcup$.
\end{itemize}
\end{definitionbox}

Now let's trace this on a concrete example, step by step, showing every tape
cell at every stage.

\begin{examplebox}{{Full Trace: Simulating a Doubly-Infinite DTM on Input ``01''}}

\textbf{Setup:} Let $M$ be a doubly-infinite tape DTM that does the following:
\begin{itemize}[nosep]
\item Starts on input $01$ at position 1.
\item Scans left from position 1 to position $-1$ (needs 3 cells to the left
  of the input).
\item Writes symbols $A$ and $B$ on the way, then accepts.
\end{itemize}

Concretely, $M$ has these transitions:
\begin{center}
\begin{tabular}{lll}
\toprule
\textbf{Current state, symbol read} & \textbf{Transition} & \textbf{Meaning} \\
\midrule
$(s, \mathtt{0})$ & $(q_1, \mathtt{0}, -1)$ & In start state, read 0: keep 0, go left \\
$(q_1, \sqcup)$ & $(q_2, A, -1)$ & Read blank: write $A$, go left \\
$(q_2, \sqcup)$ & $(t, B, +1)$ & Read blank: write $B$, accept \\
\bottomrule
\end{tabular}
\end{center}

\bigskip
\textbf{--- $M$'s behavior on the doubly-infinite tape (what we are simulating) ---}

\textbf{Step 0:} Configuration $[s, \tau, 1]$.

\begin{center}
\begin{tabular}{c|ccccccc}
Position & $-2$ & $-1$ & $0$ & $1$ & $2$ & $3$ \\
\hline
Symbol & $\sqcup$ & $\sqcup$ & $\sqcup$ & $\mathtt{0}$ & $\mathtt{1}$ & $\sqcup$ \\
\end{tabular}
\quad Head at position 1, state $s$.
\end{center}

Read $\mathtt{0}$. Transition: $\delta(s, \mathtt{0}) = (q_1, \mathtt{0}, -1)$.
Write $\mathtt{0}$ (no change), move left.

\textbf{Step 1:} Configuration $[q_1, \tau, 0]$.

\begin{center}
\begin{tabular}{c|ccccccc}
Position & $-2$ & $-1$ & $0$ & $1$ & $2$ & $3$ \\
\hline
Symbol & $\sqcup$ & $\sqcup$ & $\sqcup$ & $\mathtt{0}$ & $\mathtt{1}$ & $\sqcup$ \\
\end{tabular}
\quad Head at position 0, state $q_1$.
\end{center}

Read $\sqcup$. Transition: $\delta(q_1, \sqcup) = (q_2, A, -1)$.
Write $A$ at position 0, move left.

\textbf{Step 2:} Configuration $[q_2, \tau', -1]$.

\begin{center}
\begin{tabular}{c|ccccccc}
Position & $-2$ & $-1$ & $0$ & $1$ & $2$ & $3$ \\
\hline
Symbol & $\sqcup$ & $\sqcup$ & $A$ & $\mathtt{0}$ & $\mathtt{1}$ & $\sqcup$ \\
\end{tabular}
\quad Head at position $-1$, state $q_2$.
\end{center}

Read $\sqcup$. Transition: $\delta(q_2, \sqcup) = (t, B, +1)$.
Write $B$ at position $-1$, move right. Enter state $t$ (accept).

\textbf{Final:} Configuration $[t, \tau'', 0]$.

\begin{center}
\begin{tabular}{c|ccccccc}
Position & $-2$ & $-1$ & $0$ & $1$ & $2$ & $3$ \\
\hline
Symbol & $\sqcup$ & $B$ & $A$ & $\mathtt{0}$ & $\mathtt{1}$ & $\sqcup$ \\
\end{tabular}
\quad Head at position 0, state $t$. \textbf{ACCEPTED.}
\end{center}

$M$ used cells at positions $-1, 0, 1, 2$ --- including two cells to the left
of the input. A standard TM cannot do this directly.

\bigskip
\textbf{--- Now: $M'$ simulating $M$ on a singly-infinite tape ---}

\textbf{Initialization (Step a):} Place $\ell$ and $r$ markers around the input.

\begin{center}
\begin{tabular}{c|cccccccc}
Cell & 1 & 2 & 3 & 4 & 5 & 6 & 7 & $\cdots$ \\
\hline
Symbol & $\ell$ & $\mathtt{0}$ & $\mathtt{1}$ & $r$ & $\sqcup$ & $\sqcup$ & $\sqcup$ & $\cdots$ \\
\end{tabular}
\quad Head at cell 2, simulated state $s$.
\end{center}

The window $[\ell \cdots r]$ currently represents $M$'s positions 1 through 2.

\bigskip
\textbf{Simulating Step 0 of $M$ (Step c):}

Head reads $\mathtt{0}$ at cell 2. Look up $\delta(s, \mathtt{0}) = (q_1, \mathtt{0}, -1)$.
Write $\mathtt{0}$ (no visible change). Update simulated state to $q_1$.
Move head left to cell 1.

\begin{center}
\begin{tabular}{c|cccccccc}
Cell & 1 & 2 & 3 & 4 & 5 & 6 & 7 & $\cdots$ \\
\hline
Symbol & $\ell$ & $\mathtt{0}$ & $\mathtt{1}$ & $r$ & $\sqcup$ & $\sqcup$ & $\sqcup$ & $\cdots$ \\
\end{tabular}
\quad Head at cell 1, simulated state $q_1$.
\end{center}

\textbf{Problem!} The head is at cell 1, which contains $\ell$ --- the left
boundary marker. This means $M$ needs a cell to the left of the current
window. We must expand.

\bigskip
\textbf{Left Expansion (Step d):}

We need to shift everything from $\ell$ to $r$ one cell to the right, creating
a blank cell right after $\ell$:

\textit{Before shift:}
\begin{center}
\begin{tabular}{c|cccccccc}
Cell & 1 & 2 & 3 & 4 & 5 & 6 & 7 & $\cdots$ \\
\hline
Symbol & $\ell$ & $\mathtt{0}$ & $\mathtt{1}$ & $r$ & $\sqcup$ & $\sqcup$ & $\sqcup$ & $\cdots$ \\
\end{tabular}
\quad Head at cell 1.
\end{center}

Shift $\mathtt{0}, \mathtt{1}, r$ each one cell to the right. Insert $\sqcup$
at cell 2:

\textit{After shift:}
\begin{center}
\begin{tabular}{c|cccccccc}
Cell & 1 & 2 & 3 & 4 & 5 & 6 & 7 & $\cdots$ \\
\hline
Symbol & $\ell$ & $\sqcup$ & $\mathtt{0}$ & $\mathtt{1}$ & $r$ & $\sqcup$ & $\sqcup$ & $\cdots$ \\
\end{tabular}
\quad Head at cell 2 (the new blank cell), simulated state $q_1$.
\end{center}

The window $[\ell \cdots r]$ now represents $M$'s positions $0$ through $2$.
Cell 2 corresponds to $M$'s position 0 (the newly accessed cell, which is
blank). Cell 3 corresponds to $M$'s position 1. Cell 4 corresponds to
$M$'s position 2.

\bigskip
\textbf{Simulating Step 1 of $M$ (Step c):}

Head reads $\sqcup$ at cell 2. Look up $\delta(q_1, \sqcup) = (q_2, A, -1)$.
Write $A$ at cell 2. Update simulated state to $q_2$. Move head left to cell 1.

\begin{center}
\begin{tabular}{c|cccccccc}
Cell & 1 & 2 & 3 & 4 & 5 & 6 & 7 & $\cdots$ \\
\hline
Symbol & $\ell$ & $A$ & $\mathtt{0}$ & $\mathtt{1}$ & $r$ & $\sqcup$ & $\sqcup$ & $\cdots$ \\
\end{tabular}
\quad Head at cell 1, simulated state $q_2$.
\end{center}

\textbf{Problem again!} The head hit the $\ell$ marker \emph{again}. $M$ needs
yet another cell to the left. We expand again.

\bigskip
\textbf{Left Expansion (Step d) --- second time:}

\textit{Before shift:}
\begin{center}
\begin{tabular}{c|cccccccc}
Cell & 1 & 2 & 3 & 4 & 5 & 6 & 7 & $\cdots$ \\
\hline
Symbol & $\ell$ & $A$ & $\mathtt{0}$ & $\mathtt{1}$ & $r$ & $\sqcup$ & $\sqcup$ & $\cdots$ \\
\end{tabular}
\end{center}

Shift $A, \mathtt{0}, \mathtt{1}, r$ each one cell right. Insert $\sqcup$ at
cell 2:

\textit{After shift:}
\begin{center}
\begin{tabular}{c|cccccccc}
Cell & 1 & 2 & 3 & 4 & 5 & 6 & 7 & $\cdots$ \\
\hline
Symbol & $\ell$ & $\sqcup$ & $A$ & $\mathtt{0}$ & $\mathtt{1}$ & $r$ & $\sqcup$ & $\cdots$ \\
\end{tabular}
\quad Head at cell 2, simulated state $q_2$.
\end{center}

The window $[\ell \cdots r]$ now represents $M$'s positions $-1$ through $2$.
Cell 2 = position $-1$, cell 3 = position $0$, cell 4 = position $1$,
cell 5 = position $2$.

\bigskip
\textbf{Simulating Step 2 of $M$ (Step c):}

Head reads $\sqcup$ at cell 2. Look up $\delta(q_2, \sqcup) = (t, B, +1)$.
Write $B$ at cell 2. Update simulated state to $t$ (accept!). Move head right
to cell 3.

\begin{center}
\begin{tabular}{c|cccccccc}
Cell & 1 & 2 & 3 & 4 & 5 & 6 & 7 & $\cdots$ \\
\hline
Symbol & $\ell$ & $B$ & $A$ & $\mathtt{0}$ & $\mathtt{1}$ & $r$ & $\sqcup$ & $\cdots$ \\
\end{tabular}
\quad Head at cell 3, simulated state $t$. \textbf{ACCEPTED!}
\end{center}

\bigskip
\textbf{Verification --- do the tape contents match?}

Let's line up $M$'s final tape with $M'$'s final tape (ignoring the boundary
markers):

\begin{center}
\renewcommand{\arraystretch}{1.3}
\begin{tabular}{l|ccccc}
\toprule
\textbf{$M$'s position (doubly-infinite)} & $-1$ & $0$ & $1$ & $2$ \\
\midrule
\textbf{$M$'s final tape} & $B$ & $A$ & $\mathtt{0}$ & $\mathtt{1}$ \\
\textbf{$M'$'s tape (cells 2--5)} & $B$ & $A$ & $\mathtt{0}$ & $\mathtt{1}$ \\
\bottomrule
\end{tabular}
\end{center}

Perfect match. Both machines accepted. The simulation is correct.

\textbf{Key observations from this trace:}
\begin{itemize}[nosep]
\item $M$ took 3 steps. $M'$ took 3 simulation steps \textbf{plus} 2 shift
  operations. Each shift scans the entire window (costing $O(n)$ where $n$
  is the window size). The simulation is \textbf{slower} but \textbf{correct}.
\item The left expansion (shifting) happened twice because $M$ moved to
  positions 0 and $-1$, which were both outside the initial window. Each
  expansion grew the window by one cell.
\item The right expansion (step e) was never needed in this example because
  $M$ never moved past position 2.
\end{itemize}
\end{examplebox}

\begin{conceptbox}{Why the Simulation Is Correct}
\textbf{The Core Argument:}

$M'$ maintains the following \textbf{invariant} at all times:

\begin{quote}
The content between $\ell$ and $r$ on $M'$'s tape is an exact copy of the
corresponding region of $M$'s doubly-infinite tape. The head position on
$M'$'s tape (relative to $\ell$) corresponds to the head position on $M$'s
tape. The simulated state matches $M$'s state.
\end{quote}

\textbf{Why the invariant holds:}
\begin{enumerate}[nosep]
\item \textbf{Initially:} The input $w$ is placed between $\ell$ and $r$,
  matching $M$'s initial tape at positions $1$ through $k$. All other cells
  in $M$'s tape are blank, and our window doesn't need to include them yet.
\item \textbf{Normal simulation step:} When $M'$ simulates a transition of
  $M$, it reads the same symbol, writes the same symbol, and moves in the
  same direction. The tape content and head position stay synchronized.
\item \textbf{Left expansion:} When $M$ moves to a new cell on the left, that
  cell contains $\sqcup$ (it was never visited). The expansion creates a
  $\sqcup$ cell on $M'$'s tape in the correct position. The window grows,
  but the content still matches.
\item \textbf{Right expansion:} Symmetric argument. The new cell contains
  $\sqcup$ on both tapes.
\end{enumerate}

\textbf{Why accept/reject behavior is preserved:}

Since the invariant holds, $M'$ always reads the same symbol as $M$ at
corresponding positions. Therefore, $M'$ makes the same state transitions as
$M$. If $M$ reaches state $t$, so does $M'$. If $M$ reaches state $r$, so
does $M'$. If $M$ loops, $M'$ loops too (possibly slower, but forever).

\kt{The simulation works because $M'$ maintains a perfect copy of $M$'s
active tape region between the boundary markers $\ell$ and $r$. Whenever $M$
needs a new cell, $M'$ expands the window. The shifting is expensive --- each
left expansion costs $O(n)$ steps --- but it is always finite, and the
content always stays synchronized.}
\end{conceptbox}


%───────────────────────────────────────────────────────────────────────────────
\subsection{Part 3: Doubly-Infinite Simulates Singly-Infinite (The Easy Direction)}
\label{subsec:doubly-simulates-singly}
%───────────────────────────────────────────────────────────────────────────────

\begin{notebox}{Analogy: A Bigger Room Contains a Smaller Room}
Imagine you have a large warehouse with infinite space in both directions
(left and right). Someone gives you a floor plan that only uses the right
half of the warehouse (rooms 1, 2, 3, \ldots). Can you implement that floor
plan in your warehouse?

Of course! You just use the right half and ignore the left half. There's only
one catch: the floor plan says ``if you walk into the left wall at room 1,
you bounce off and stay in room 1.'' Your warehouse doesn't have a wall
there --- room 0 exists. So you \textbf{build a wall}: put a barrier at
position 0 that forces anyone who hits it to bounce back to position 1.

\textbf{In our formal world:}
\begin{itemize}[nosep]
\item Large warehouse = doubly-infinite tape (positions $\mathbb{Z}$)
\item Right half = cells $1, 2, 3, \ldots$ (same as a singly-infinite tape)
\item Left wall at room 1 = the boundary behavior of a singly-infinite tape
  (head at position 1 trying to go left stays at 1)
\item Barrier at position 0 = a special marker symbol \texttt{\#} placed at
  position 0 that causes the head to ``bounce right''
\end{itemize}
\end{notebox}

\begin{definitionbox}{{Simulation Construction: Doubly-Infinite Simulates Singly-Infinite}}
\textbf{Plain English:}
Given a singly-infinite tape DTM $M$, we build a doubly-infinite tape DTM
$M'$ that does the same thing. The cells $1, 2, 3, \ldots$ are used normally.
To simulate the left wall, $M'$ places a marker \texttt{\#} at position 0.
Whenever the head hits \texttt{\#}, it ``bounces'' to the right --- mimicking
the singly-infinite behavior of ``head stays at position 1 when trying to
move left.''

\textbf{Formal construction:}

Let $M = (Q, \Sigma, \Gamma, s, t, r, \delta)$ be a singly-infinite tape DTM.
Define $M' = (Q', \Sigma, \Gamma', s', t, r, \delta')$ as a doubly-infinite
tape DTM where:

\begin{itemize}[nosep]
\item $Q' = Q \cup \{s', q_\#\}$ where $s', q_\#$ are fresh states not in $Q$
\item $\Gamma' = \Gamma \cup \{\#\}$ where $\# \notin \Gamma$ is a fresh marker
  symbol
\item $s'$ is the new start state (not $s$ --- we need two setup steps first)
\end{itemize}

\textbf{Transition function $\delta'$:}

\textbf{(1) Initialization --- place the marker:}
\begin{itemize}[nosep]
\item $\delta'(s', a) = (q_\#, a, -1)$ for all $a \in \Gamma'$

  \textit{``In the new start state $s'$, no matter what symbol we read,
  keep that symbol, and move left to position 0.''}
\item $\delta'(q_\#, \sqcup) = (s, \#, +1)$

  \textit{``At position 0 (which is blank), write the marker \texttt{\#},
  then move right back to position 1 and enter $M$'s original start state
  $s$.''}
\end{itemize}

After these two transitions, the tape looks like:
\begin{lstlisting}
... | _ | # | w1 | w2 | ... | wk | _ | ...
         0    1     2          k
              ^
        head here, state s
\end{lstlisting}
Position 0 has the marker \texttt{\#}. The head is at position 1, in state
$s$ --- exactly how $M$ would start.

\textbf{(2) Normal simulation:}
\begin{itemize}[nosep]
\item $\delta'(q, a) = \delta(q, a)$ for all $q \in Q \setminus \{t, r\}$
  and all $a \in \Gamma$

  \textit{``For any non-halting state of $M$ and any symbol from $M$'s
  alphabet, do exactly what $M$ would do.''}
\end{itemize}

\textbf{(3) Boundary bounce --- the wall simulation:}
\begin{itemize}[nosep]
\item $\delta'(q, \#) = (q, \#, +1)$ for all $q \in Q \setminus \{t, r\}$

  \textit{``If any non-halting state reads the \texttt{\#} marker, don't
  change state, keep the marker, and move right. This simulates the
  singly-infinite tape behavior where the head bounces off the left wall.''}
\end{itemize}

\textbf{Breaking Down the Construction:}
\begin{itemize}[nosep]
\item \textbf{$s'$ and $q_\#$:} Two new states used only during initialization.
  Once the marker is placed, these states are never visited again.
\item \textbf{The marker \texttt{\#}:} Sits at position 0 for the entire
  computation. It never gets overwritten (rule 3 always writes \texttt{\#}
  back). It acts as an artificial left wall.
\item \textbf{The bounce (rule 3):} On a singly-infinite tape, if $M$ is in
  state $q$ at position 1 and tries to move left, the head stays at position
  1 and $M$ continues in whatever state the transition specifies. On our
  doubly-infinite simulation, the head would move to position 0 and read
  \texttt{\#}. Rule 3 sends the head back to position 1 \emph{without
  changing state}. The net effect is the same: the head ``stays'' at
  position 1.
\item \textbf{Why we don't change state in the bounce:} On a singly-infinite
  tape, when the head ``stays at position 1,'' it still transitions to the
  new state specified by $\delta$. That state change already happened in
  step (2) when the transition was executed. Rule (3) just handles the
  physical bounce-back --- no additional state change needed.
\end{itemize}
\end{definitionbox}

\begin{examplebox}{{Trace: Doubly-Infinite Simulating a Singly-Infinite DTM (Bounce Example)}}

\textbf{Setup:} Let $M$ be a singly-infinite tape DTM with these transitions:
\begin{center}
\begin{tabular}{lll}
\toprule
\textbf{State, symbol} & \textbf{Transition} & \textbf{Meaning} \\
\midrule
$(s, \mathtt{0})$ & $(q_1, X, -1)$ & Read 0: write $X$, go left \\
$(q_1, \mathtt{0})$ & $(t, \mathtt{0}, +1)$ & Read 0: accept \\
\bottomrule
\end{tabular}
\end{center}

On input $\mathtt{01}$ with a \textbf{singly-infinite} tape:
\begin{enumerate}[nosep]
\item $[s, \text{tape}, 1]$: Read $\mathtt{0}$. Write $X$, go left. New state $q_1$.
  But head is at position 1 and tries to go left $\to$ \textbf{stays at position 1} (left wall!).
\item $[q_1, \text{tape}, 1]$: Read $X$ (we just wrote it). Hmm --- $\delta(q_1, X)$ is undefined.
\end{enumerate}

Wait --- that doesn't work well. Let's pick a cleaner example. Suppose $M$ has:
\begin{center}
\begin{tabular}{lll}
\toprule
\textbf{State, symbol} & \textbf{Transition} & \textbf{Meaning} \\
\midrule
$(s, \mathtt{0})$ & $(q_1, \mathtt{0}, -1)$ & Read 0: keep 0, go left \\
\bottomrule
\end{tabular}
\end{center}

On input $\mathtt{01}$, singly-infinite tape:
\begin{enumerate}[nosep]
\item $[s, \text{tape}, 1]$: Read $\mathtt{0}$. $\delta(s, 0) = (q_1, 0, -1)$.
  Write 0 (no change), try to move left. Position is 1 $\to$ head \textbf{stays at 1}. State becomes $q_1$.
\item $[q_1, \text{tape}, 1]$: Head is at position 1, reads $\mathtt{0}$.
  (The computation continues from here with whatever $\delta(q_1, 0)$ says.)
\end{enumerate}

The key point: the head bounced off the left wall. It tried to go to position 0
but couldn't.

\textbf{Now $M'$ (doubly-infinite) simulating $M$:}

\textbf{Initialization:}

Start: head at position 1, state $s'$.

\begin{center}
\begin{tabular}{c|ccccccc}
Position & $-1$ & $0$ & $1$ & $2$ & $3$ \\
\hline
Symbol & $\sqcup$ & $\sqcup$ & $\mathtt{0}$ & $\mathtt{1}$ & $\sqcup$ \\
\end{tabular}
\quad Head at 1, state $s'$.
\end{center}

$\delta'(s', \mathtt{0}) = (q_\#, \mathtt{0}, -1)$: keep the 0, move left.
Head goes to position 0.

\begin{center}
\begin{tabular}{c|ccccccc}
Position & $-1$ & $0$ & $1$ & $2$ & $3$ \\
\hline
Symbol & $\sqcup$ & $\sqcup$ & $\mathtt{0}$ & $\mathtt{1}$ & $\sqcup$ \\
\end{tabular}
\quad Head at 0, state $q_\#$.
\end{center}

$\delta'(q_\#, \sqcup) = (s, \#, +1)$: write \texttt{\#}, move right.

\begin{center}
\begin{tabular}{c|ccccccc}
Position & $-1$ & $0$ & $1$ & $2$ & $3$ \\
\hline
Symbol & $\sqcup$ & $\#$ & $\mathtt{0}$ & $\mathtt{1}$ & $\sqcup$ \\
\end{tabular}
\quad Head at 1, state $s$. Initialization complete.
\end{center}

\textbf{Simulating Step 1 of $M$:}

$\delta'(s, \mathtt{0}) = \delta(s, \mathtt{0}) = (q_1, \mathtt{0}, -1)$:
write 0 (no change), move left. Head goes to position \textbf{0}.

\begin{center}
\begin{tabular}{c|ccccccc}
Position & $-1$ & $0$ & $1$ & $2$ & $3$ \\
\hline
Symbol & $\sqcup$ & $\#$ & $\mathtt{0}$ & $\mathtt{1}$ & $\sqcup$ \\
\end{tabular}
\quad Head at 0, state $q_1$.
\end{center}

The head reads \texttt{\#}! This triggers the \textbf{bounce rule}:
$\delta'(q_1, \#) = (q_1, \#, +1)$: keep state $q_1$, keep \texttt{\#},
move right.

\begin{center}
\begin{tabular}{c|ccccccc}
Position & $-1$ & $0$ & $1$ & $2$ & $3$ \\
\hline
Symbol & $\sqcup$ & $\#$ & $\mathtt{0}$ & $\mathtt{1}$ & $\sqcup$ \\
\end{tabular}
\quad Head at 1, state $q_1$. \textbf{Bounce complete!}
\end{center}

\textbf{Result:} After the bounce, the head is at position 1 in state $q_1$
--- exactly where it would be on the singly-infinite tape after bouncing off
the left wall.

\kt{The bounce works in two micro-steps on $M'$ (move to \# at position 0,
then bounce back to position 1) but produces the same result as one step on
$M$ (head stays at position 1). The simulation is correct because the net
effect on state and head position is identical.}
\end{examplebox}


%───────────────────────────────────────────────────────────────────────────────
\subsection{The Conclusion: Doubly-Infinite = Singly-Infinite}
\label{subsec:doubly-infinite-conclusion}
%───────────────────────────────────────────────────────────────────────────────

\begin{conceptbox}{Putting Both Directions Together}
We have shown:

\textbf{Direction 1 (Part 2):} Every singly-infinite tape DTM can simulate
every doubly-infinite tape DTM. Given any doubly-infinite machine $M$, we
built a singly-infinite machine $M'$ with boundary markers $\ell$ and $r$
that tracks $M$'s active region and expands as needed.
\begin{center}
Doubly-infinite tape DTMs $\subseteq$ Singly-infinite tape DTMs (in power)
\end{center}

\textbf{Direction 2 (Part 3):} Every doubly-infinite tape DTM can simulate
every singly-infinite tape DTM. Given any singly-infinite machine $M$, we
built a doubly-infinite machine $M'$ that places a marker \texttt{\#} at
position 0 and bounces the head whenever it hits the marker.
\begin{center}
Singly-infinite tape DTMs $\subseteq$ Doubly-infinite tape DTMs (in power)
\end{center}

\textbf{Both directions together:}
\fitmath{\text{Languages decided by singly-infinite DTMs} \;=\; \text{Languages decided by doubly-infinite DTMs}}

The doubly-infinite tape is a convenience, not a power upgrade. You can always
convert between the two models. The class of decidable (and c.e.) languages
remains the same.
\end{conceptbox}

\begin{conceptbox}{{The Robustness Pattern --- Four Proofs, One Template}}
Looking back at Sections 3 through 6, every robustness proof followed the
\textbf{same three-step template}:

\begin{center}
\renewcommand{\arraystretch}{1.4}
\begin{tabular}{clll}
\toprule
\textbf{Sec.} & \textbf{Extension} & \textbf{Simulation Trick} & \textbf{Cost} \\
\midrule
3 & Stay direction & Encode ``stay'' as right-then-left & 2 extra steps per stay \\
4 & Multi-tape & Encode $k$ tapes on 1 tape & $O(n)$ overhead per step \\
5 & Nondeterminism & BFS over computation tree & Exponential blowup \\
6 & Doubly-infinite tape & Boundary markers + shifting & $O(n)$ per left expansion \\
\bottomrule
\end{tabular}
\end{center}

The template:
\begin{enumerate}[nosep]
\item \textbf{Define the extension:} What does the fancy machine look like formally?
\item \textbf{Build a simulation:} Construct a standard DTM that mimics the fancy machine step by step.
\item \textbf{Argue correctness:} Show the simulator produces the same accept/reject/loop on every input.
\end{enumerate}

Each simulation is slower than the original (sometimes dramatically slower),
but speed doesn't matter for robustness --- only the final verdict matters.

\kt{Together with stay direction (Section 3), multi-tape (Section 4), and
nondeterminism (Section 5), the doubly-infinite tape gives us four robustness
proofs. Each follows the same template: define the extension, build a
simulation, trace it, argue correctness. The TM model is remarkably robust
--- no ``reasonable'' extension adds computational power. This is strong
evidence for the Church-Turing Thesis.}
\end{conceptbox}
